\subsection{Latent dynamics and learning objectives}
\label{sec:C06:latent}
Let $z_t$ be a stochastic latent variable and $h_t$ a deterministic recurrent state. A generic latent dynamics model uses a prior $p_\phi(z_t\mid h_t)$, an observation-conditioned posterior $q_\phi(z_t\mid h_t,o_t)$, and an update $h_{t+1}=f_\phi(h_t,z_t,a_t)$. Observation, reward, and termination heads provide learning signals. A variational training objective combines reconstruction or prediction terms with a divergence between posterior and prior. Its purpose is to make a state inferable from real observations while allowing future latent trajectories to be sampled without access to those future observations.

The two directions of information flow matter. The posterior can use the new observation to correct its state estimate. The prior must predict before seeing that observation. A training loss that allows the prior to condition on future information is not a valid test of imagination. Similarly, a replay dataset must distinguish what was known at action time from what was later discovered. The same temporal leakage that damages agent policy training can make a world model appear more predictive than its deployable interface permits. Saving full trajectories is not enough unless their causal visibility is retained.

DreamerV3 provides a concrete latent-world-model approach in which the actor and critic learn from imagined latent trajectories. Its reported method uses separate prediction, dynamics, and representation losses, with stop-gradient placement and regularization controlling their interaction. The deployed actor in the cited version samples actions without online lookahead planning~\cite{hafner2023dreamerv3}. This distinction prevents a common category error: learning through imagined trajectories and searching a model at every decision are different uses of a world model. Both can improve control, but their computation and failure modes differ.

Representation pressure creates a tradeoff. Strong compression can make dynamics easier to predict while discarding task-critical detail. Weak compression can preserve observations while making the latent transition harder to learn. Reward and continuation predictions encourage retention of some control-relevant information, but their sufficiency depends on the training objectives and task family. A reward change can reveal distinctions the original representation ignored. A world model intended for reusable enterprise agents should therefore be evaluated under changed goals and tools, not only the reward function that shaped its latent state.

An uncertainty estimate should distinguish ambiguity in observations, stochastic environmental outcomes, and lack of model knowledge. Sampling multiple futures from one generative model does not automatically capture the third category. All samples can share the same systematic error. Ensembles, held-out transition tests, and disagreement across independently informed predictors can provide diagnostics, but dependence and shared training data limit their interpretation. The controller needs a policy for falling back to real observations or conservative actions when model evidence is weak; otherwise uncertainty remains an unused output rather than a control mechanism.

\subsection{Two temporal scales require an executable refinement interface}
\label{sec:C06:dualscale}
Dual-WM's September 2026 construction assigns primitive transitions and macro-action transitions to distinct latent spaces. Its Section 3 encodes windows of low-level features into a high-level state, learns a regularized macro-action interface, and refines high-level subgoals by projecting low-level rollouts back into the high-level space~\cite{rewrite_dualwm2026}. It does not require an inverse decoder from an arbitrary high-level state to one concrete primitive state. A final low-level goal phase restores distinctions that an abstract subgoal can omit.

The temporal units matter. If one low-level model action spans $f$ environment operations and one macro transition spans $k$ low-level steps, a high-level horizon $H_H$ covers $kfH_H$ operations. Comparing $H_H$ directly with a primitive planning horizon confounds abstraction with temporal reach. Moreover, a high-level transition is usable only if some allowed primitive sequence realizes its proposed subgoal. A normal prior over macro-actions regularizes search but does not prove that every sampled point has an executable refinement.

For an agent experiment, fix an overall operation budget and compare primitive planning, two-scale planning without refinement validation, and two-scale planning with observed subgoal checks. Hold the number of environment operations covered constant before comparing prediction depth or latency. Include abstract states that alias different permissions or irreversible changes. If refinement reaches the high-level embedding but violates a required concrete predicate, score it as a failed refinement. This separates useful temporal abstraction from a planner that merely hides execution error inside a coarser representation.
